\subsection{EXTENSION OF A CONTINUOUS REAL-VALUED FUNCTION}

Let X and Y be two topological spaces and let \(A \neq X\) be a closed subset of X. If f is a continuous mapping of A into Y, it is not always possible to extend f to a continuous mapping of the whole of X into Y. When \(Y = \overline{R}\) , the possibility of such an extension is determined by the following theorem:

THEOREM 2 (Urysohn). Axiom \((\mathbf{O}_{\mathbf{V}})\) is equivalent to the following property: \((\mathbf{O}_{\mathbf{V}}^{\prime \prime \prime})\) Given any closed subset A of X and any continuous real-valued function \(f\) (finite or not) defined on A, there exists an extension \(g\) of \(f\) to the whole space X, which is a continuous mapping of X into \(\overline{\mathbf{R}}\) .

It is easy to see that \((\mathbf{O}_{\mathbf{V}}^{\prime\prime})\) implies \((\mathbf{O}_{\mathbf{V}})\) ; for if B and C are two disjoint closed subsets of X, then the function which is equal to o on B and equal to \(\mathbf{I}\) on \(\mathbf{C}\) is defined and continuous on the closed set \(\mathbf{B} \cup \mathbf{C}\) , hence by \((\mathrm{O}_{\mathbf{V}}^{\prime \prime})\) has a continuous extension \(f\) to \(\mathbf{X}\) . If \(g = \inf (f^{+}, \mathbf{I})\) , then \(g\) is continuous on \(\mathbf{X}\) , takes its values in \([0, \mathbf{I}]\) and is equal to \(\mathbf{0}\) on \(\mathbf{B}\) and to \(\mathbf{I}\) on \(\mathbf{C}\) .

Let us show conversely that \((\mathbf{O}_{\mathbf{v}})\) implies \((\mathbf{O}_{\mathbf{v}}^{\prime\prime})\) . Since \(\overline{R}\) and the interval \([-1, +1]\) are homeomorphic, we need consider only the case where the continuous mapping \(f: A \to \overline{R}\) takes its values in \([-1, +1]\) . We shall construct an extension g of f to X by forming a sequence \((g_{n})\) of continuous functions on X, such that the sequence \((g_{n}(x))\) converges for all \(x \in X\) to a point of the interval \([-1, +1]\) ; this limit will, by definition, be the value of g at x, and it will follow from the choice of the \(g_{n}\) that the function g satisfies the required conditions.

The definition of the \(g_{n}\) rests on the following lemma:

LEMMA 1. Let u be a continuous mapping of A into \([-1, +1]\) ; then there is a continuous mapping v of X into \([-1/3, +1/3]\) , such that \(|u(x) - v(x)| \leqslant 2/3\) for all \(x \in A\) .

Let H be the set of all \(x \in A\) such that \(-\mathtt{I} \leqslant u(x) \leqslant -\mathtt{I}/3\) , and let K be the set of all \(x \in A\) such that \(\mathtt{I}/3 \leqslant u(x) \leqslant \mathtt{I}\) ; H and K are closed in A, and therefore in X, and do not intersect; hence by \((\mathrm{O}_{\mathrm{V}})\) there is a continuous mapping v of X into \([-\mathtt{I}/3, +\mathtt{I}/3]\) which is equal to \(-\mathtt{I}/3\) on H and to \(+\mathtt{I}/3\) on K. The mapping v satisfies the conditions of the lemma.

We now define the functions \(g_n\) by induction. Applying the lemma with \(u = f\) , we define \(g_0\) to be a continuous mapping of \(\mathbf{X}\) into \([-1/3, +1/3]\) such that \(|f(x) - g_0(x)| \leqslant 2/3\) for all \(x \in \mathbf{A}\) . Suppose now that a continuous mapping \(g_n\) of \(\mathbf{X}\) into the interval

\[
[ - \mathrm{I} + (\frac {2}{3}) ^ {n + 1}, \mathrm{I} - (\frac {2}{3}) ^ {n + 1} ]
\]

has been defined, such that \(|f(x) - g_n(x)| \leqslant (\frac{2}{3})^{n+1}\) for all \(x \in \mathbf{A}\) . Applying the lemma to the function \(u(x) = (\frac{2}{3})^{n+1} (f(x) - g_n(x))\) , we see that there exists a continuous mapping \(h_{n+1}\) of \(\mathbf{X}\) into the interval \(\left[-\frac{2^{n+1}}{3^{n+2}}, \frac{2^{n+1}}{3^{n+2}}\right]\) such that

\[
\left| f (x) - g _ {n} (x) - h _ {n + 1} (x) \right| \leqslant (2 / 3) ^ {n + 2}
\]

for all \(x \in \mathbf{A}\) ; the induction is completed by taking \(g_{n+1} = g_n + h_{n+1}\) , since this function satisfies the inequality \(|g_{n+1}(x)| \leqslant 1 - (2/3)^{n+2}\) for all \(x \in \mathbf{X}\) , by virtue of the definition of \(h_{n+1}\) .

From this definition it follows that, for \(m \geqslant p\) and \(n \geqslant p\) , we have

\[
\left| g _ {m} (x) - g _ {n} (x) \right| \leqslant \frac {2 ^ {p + 1}}{3 ^ {p + 2}} \sum_ {k = 0} ^ {\infty} (2 / 3) ^ {k} = (2 / 3) ^ {p + 1}
\]

at each point \(x \in \mathbf{X}\) ; hence the sequence \((g_n(x))\) is a Cauchy sequence for each \(x \in \mathbf{X}\) , and therefore converges to a point \(g(x)\) of the interval \([-1, +1]\) ; and since \(f(x) - g_n(x)\) tends to o for all \(x \in \mathbf{A}\) as \(n \to \infty\) , \(g\) is an extension of \(f\) to \(\mathbf{X}\) . It remains therefore only to show that \(g\) is continuous on \(\mathbf{X}\) .

Now let \(x\) be any point of \(\mathbf{X}\) ; then, given any \(\varepsilon > 0\) , there exists an integer \(n_0\) such that \(|g_m(y) - g_n(y)| \leqslant \varepsilon\) for all \(y \in \mathbf{X}\) and all \(m \geqslant n_0\) and all \(n \geqslant n_0\) ; hence, letting \(m\) tend to \(+\infty\) , we have

\[
\left| g (y) - g _ {n} (y) \right| \leqslant \varepsilon .
\]

Let \(\mathbf{V}\) be a neighbourhood of \(x\) such that \(|g_n(y) - g_n(x)| \leqslant \varepsilon\) for all \(y \in \mathbf{V}\) ; then, for each \(y \in \mathbf{V}\) we shall have

\[
\left| g (y) - g (x) \right| \leqslant \left| g (y) - g _ {n} (y) \right| + \left| g _ {n} (y) - g _ {n} (x) \right| + \left| g (x) - g _ {n} (x) \right| \leqslant 3 \varepsilon ,
\]

which shows that g is continuous at x, and completes the proof of Theorem 2. (The last part of the proof uses, in a particular case, the idea of uniform convergence, which we shall study in general in Chapter X, no. 1.)

COROLLARY. If \(f\) is a finite continuous real-valued function defined on \(\mathbf{A}\) , then there exists a finite continuous real-valued function \(g\) defined on \(\mathbf{X}\) , which extends \(f\) .

First consider the case in which \(f(x) \geqslant 0\) for all \(x \in A\) ; then there is a continuous extension \(g_1\) of \(f\) to \(X\) , taking its values in \([0, +\infty]\) . If we put \(B = \overline{g_1} (+\infty)\) , then \(B\) is closed and by hypothesis does not meet \(A\) ; the function \(h\) which is equal to \(f\) on \(A\) and to \(o\) on \(B\) is therefore a continuous function on the closed set \(A \cup B\) . Let \(g_2\) be a continuous extension of \(h\) to \(X\) , again taking its values in \([0, +\infty]\) ; the function \(g = \inf(g_1, g_2)\) is then a continuous extension of \(f\) to \(X\) , whose values are \(\geqslant 0\) and finite at every point of \(X\) .

To pass to the general case, it is enough to remark that, if \(f\) is finite and continuous on \(\mathbf{A}\) , then so are \(f^{+}\) and \(f^{-}\) ; extending \(f^{+}\) and \(f^{-}\) to \(\mathbf{X}\) by finite continuous functions \(g_{1}\) and \(g_{2}\) respectively, we see that the function \(g_{1} - g_{2}\) is finite and continuous on \(\mathbf{X}\) and extends \(f\) .

Remark. If X is a normal space and if A is a closed subset of X, there exists also a continuous extension to X of every continuous mapping f of A into a cube \(K^{I}\) (§ I, no. 5); for we have then \(f = (f_{i})_{i \in I}\) , \(f_{i}\) being a continuous mapping of A into the compact interval K of R; since there exists a continuous mapping \(g_{i}: X \to K\) which extends \(f_{i}\) , the mapping \(g = (g_{i})\) is a continuous extension of f to X.

